\section{USAMO 1979}
        \begin{enumerate}
        	\item (\textit{USAMO 1979}) Determine all non-negative integral solutions $(n_1,n_2,\dots , n_{14})$ if any, apart from permutations, of the Diophantine Equation $n_1^4+n_2^4+\cdots +n_{14}^4=1599$.


 



\textbf{Solution  1}

Recall that $n_i^4\equiv 0,1\bmod{16}$ for all integers $n_i$. Thus the sum we have is anything from 0 to 14 modulo 16. But $1599\equiv 15\bmod{16}$, and thus there are no integral solutions to the given Diophantine equation.

 



\textbf{Solution  2}

In base $16$, this equation would look like:
 \[n_1^4+n_2^4+\cdots +n_{14}^4=63F_{16}\]

 We notice that the unit digits of the LHS of this equation should equal to $F_{16}$. In base $16$, the only unit digits of fourth powers are $0$ and $1$. Thus, the maximum of these $14$ terms is 14 $1's$ or $E_{16}$. Since $E_{16}$ is less than $F_{16}$, there are no integral solutions for this equation.

 


	\item (\textit{USAMO 1979}) $N$ is the north pole. $A$ and $B$ are points on a great circle through $N$ equidistant from $N$. $C$ is a point on the equator. Show that the great circle through $C$ and $N$ bisects the angle $ACB$ in the spherical triangle $ABC$ (a spherical triangle has great circle arcs as sides).


 



\textbf{Solution}


                    \begin{figure}[H]
                        \centering
                        \includegraphics[width=0.5\textwidth]{Solution/USAMO_I_1979_2_Solution_Figure_1.png}
                    \end{figure}
                    

 Since $N$ is the north pole, we define the Earth with a sphere of radius one in space with $N=(0,0,1)$ and sphere center $O=(0,0,0)$We then pick point $N$ on the sphere and define the $xz$-plane as the plane that contains great circle points $A$ , $B$, and $N$ with the $x$-axis perpendicular to the $z$-axis and in the direction of $A$.

 Using this coordinate system and $x$, $y$, and $z$ axes $A=(cos(\phi),0,sin(\phi))$ where $\phi$ is the angle from the $xy$-plane to $A$ or latitude on this sphere with $\frac{-\pi}{2} < \phi < \frac{\pi}{2}$

 Since $A$ and $B$ are points on a great circle through $N$ equidistant from $N$, then $B=(-cos(\phi),0,sin(\phi))$ 

 Since $C$ is a point on the equator, then $C=(cos(\theta),sin(\theta),0)$ where $\theta$ is the angle on the $xy$-plane from the origin to $C$ or longitude on this sphere with $-\pi < \phi \le \pi$

 We note that vectors from the origin to points $N$, $A$, $B$, and $C$ are all unit vectors because all those points are on the unit sphere.

 So, we're going to define points $N$, $A$, $B$, and $C$ as unit vectors with their coordinates.

 We also define the following vectors as follows:

 Vector $\overrightarrow{V_{CN}}$ is the unit vector in the direction of arc $CN$ and tangent to the great circle of $CN$ at $C$

 Vector $\overrightarrow{V_{CA}}$ is the unit vector in the direction of arc $CA$ and tangent to the great circle of $CA$ at $C$

 Vector $\overrightarrow{V_{CB}}$ is the unit vector in the direction of arc $CB$ and tangent to the great circle of $CB$ at $C$

 To calculate each of these vectors we shall use the cross product as follows:

 $\overrightarrow{V_{CN}}=(\overrightarrow{C}\times\overrightarrow{N})\times\overrightarrow{C}$

 $\overrightarrow{V_{CN}}=(\left\langle cos(\theta),sin(\theta),0 \right\rangle\times\left\langle 0,0,1 \right\rangle)\times \left\langle cos(\theta),sin(\theta),0 \right\rangle$

 $\overrightarrow{V_{CN}}=\left\langle sin(\theta),-cos(\theta),0 \right\rangle\times \left\langle cos(\theta),sin(\theta),0 \right\rangle$

 $\overrightarrow{V_{CN}}=\left\langle 0,0,sin^{2}(\theta)+cos^{2}(\theta) \right\rangle$

 $\overrightarrow{V_{CN}}=\left\langle 0,0,1 \right\rangle$

 Vector $\overrightarrow{V_{CA}}$:

 $\overrightarrow{V_{CA}}=(\overrightarrow{C}\times\overrightarrow{A})\times\overrightarrow{C}$

 $\overrightarrow{V_{CA}}=(\left\langle cos(\theta),sin(\theta),0 \right\rangle\times\left\langle cos(\phi),0,sin(\phi) \right\rangle)\times \left\langle cos(\theta),sin(\theta),0 \right\rangle$

 $\overrightarrow{V_{CA}}=\left\langle sin(\theta)sin(\phi),-cos(\theta)sin(\phi),-sin(\theta)cos(\phi) \right\rangle\times \left\langle cos(\theta),sin(\theta),0 \right\rangle$

 Since we're only interested in the $z$ component of the vector

 $\overrightarrow{V_{CA}}=\left\langle V_{CA_{x}},V_{CA_{y}},sin^{2}(\theta)sin(\phi)+cos^{2}(\theta)sin(\phi) \right\rangle$

 $\overrightarrow{V_{CA}}=\left\langle V_{CA_{x}},V_{CA_{y}},sin(\phi) \right\rangle$

 Vector $\overrightarrow{V_{CB}}$:

 $\overrightarrow{V_{CB}}=(\overrightarrow{C}\times\overrightarrow{b})\times\overrightarrow{C}$

 $\overrightarrow{V_{CB}}=(\left\langle cos(\theta),sin(\theta),0 \right\rangle\times\left\langle -cos(\phi),0,sin(\phi) \right\rangle)\times \left\langle cos(\theta),sin(\theta),0 \right\rangle$

 $\overrightarrow{V_{CB}}=\left\langle sin(\theta)sin(\phi),-cos(\theta)sin(\phi),sin(\theta)cos(\phi) \right\rangle\times \left\langle cos(\theta),sin(\theta),0 \right\rangle$

 Since we're only interested in the $z$ component of the vector

 $\overrightarrow{V_{CB}}=\left\langle V_{CB_{x}},V_{CB_{y}},sin^{2}(\theta)sin(\phi)+cos^{2}(\theta)sin(\phi) \right\rangle$

 $\overrightarrow{V_{CB}}=\left\langle V_{CB_{x}},V_{CB_{y}},sin(\phi) \right\rangle$

 Since we're working with unit vectors, then we can use dot products on the vectors with their angles as follows:

 $cos(\angle ACN) = \overrightarrow{V_{CA}}\cdot \overrightarrow{V_{CN}}$

 $cos(\angle ACN) = \left\langle V_{CA_{x}},V_{CA_{y}},sin(\phi) \right\rangle \cdot \left\langle 0,0,1 \right\rangle = 0*V_{CA_{x}}+0*V_{CA_{y}}+1*sin(\phi)=sin(\phi)$

 Likewise,

 $cos(\angle BCN) = \overrightarrow{V_{CB}}\cdot \overrightarrow{V_{CN}}$

 $cos(\angle BCN) = \left\langle V_{CB_{x}},V_{CB_{y}},sin(\phi) \right\rangle \cdot \left\langle 0,0,1 \right\rangle = 0*V_{CB_{x}}+0*V_{CB_{y}}+1*sin(\phi)=sin(\phi)$

 Therefore,

 $cos(\angle ACN) = cos(\angle BCN)$ and thus $\angle ACN = \angle BCN$

 Since those angles are equal, it proves that the great circle through $C$ and $N$ bisects the $\angle ACB$ in the spherical triangle $ABC$ 

 <i>Alternate solutions are always welcome. If you have a different, elegant solution to this problem, please \href{https://artofproblemsolving.com/wiki/index.php?title=1979\_USAMO\_Problems/Problem\_2&action=edit}{add it} to this page.</i>

 


	\item (\textit{USAMO 1979}) $a_1, a_2, \ldots, a_n$ is an arbitrary sequence of positive integers. A member of the sequence is picked at 
random. Its value is $a$. Another member is picked at random, independently of the first. Its value is $b$. Then a third value, $c$. Show that the probability that $a + b + c$ is divisible by $3$ is at least $\frac14$.


 



\textbf{Solution}

Let x equal the probability of picking an element of the sequence equivalent to 0 mod 3, y equal those 1 mod 3, and z equal those 2 mod 3. Then, considering that 0+0+0, 1+1+1, 2+2+2, and 0+1+2 are divisible by 3, we obtain the equivalent inequality in the First Hint. This simplifies to $x^3 + y^3 + z^3 + 6xyz \ge x^2(y+z) + y^2(x+z) + z^2(x+y).$ Now, we attempt to prove that $x^3 + y^3 + z^3 + 3xyz \ge x^2(y+z) + y^2(x+z) + z^2(x+y).$ for all nonnegative real numbers x, y, z; our inequality will follow from this. But this is just Schur's Inequality with r = 1! This completes the proof.

 \textit{Note: Schur's Inequality states that $x^r(x-y)(x-z) + y^r(y-z)(y-x) + z^r(z-x)(z-y) \ge 0$. This can be proven by setting WLOG $x \ge y \ge z$ and bunching together terms that are nonnegative. Challenge: when does equality hold? When does equality hold in the problem?}

 


	\item (\textit{USAMO 1979}) $P$ lies between the rays $OA$ and $OB$. Find $Q$ on $OA$ and $R$ on $OB$ collinear with $P$ so that $\frac{1}{PQ} + \frac{1}{PR}$ is as large as possible.


 



\textbf{Solution (inversions)}

Perform the inversion with center $P$ and radius $\overline{PO}.$ Lines $OA,OB$ go to the circles $(O_1),(O_2)$ passing through $P,O$ and the line $QR$ cuts $(O_1),(O_2)$ again at the inverses $Q',R'$ of $Q,R.$ Hence

 $\frac{1}{PQ}+\frac{1}{PR}=\frac{PQ'+PR'}{PO^2}=\frac{Q'R'}{PO^2}$

 Thus, it suffices to find the line through $P$ that maximizes the length of the  segment $\overline{Q'R'}.$ If $M,N$ are the midpoints of $PQ',PR',$ i.e. the projections of $O_1,O_2$ onto $QR,$ then from the right trapezoid $O_1O_2NM,$ we deduce that $O_1O_2 \ge MN = \frac{_1}{^2}Q'R'.$ Consequently, $2 \cdot O_1O_2$ is the greatest possible length of $Q'R',$ which obviously occurs when $O_1O_2NM$ is a rectangle. Hence, $Q,R$ are the intersections of $OA,OB$ with the perpendicular to $PO$ at $P.$

 



\textbf{Solution (trig bash)}


                    \begin{figure}[H]
                        \centering
                        \includegraphics[width=0.5\textwidth]{Solution/USAMO_I_1979_4_Solution_Figure_1.png}
                    \end{figure}
                    

 Let $r = OP, x = \angle OPR, a = \angle POR,$ and $b = \angle POQ.$ Then $\angle ORP = \pi - x - a$ and $\angle OQP = x - b.$ Using the Law of Sines on $\triangle OPR$ gives
 \[PR = \sin a * \frac{r}{\sin(\pi - x - a)} = \sin a * \frac{r}{\sin(x + a)},\]and using the Law of Sines on $\triangle OPQ$ gives
 \[PQ = \sin b * \frac{r}{\sin(x - b)}.\]Note that $r, a,$ and $b$ are given constants.Hence,
 \[\frac{1}{PR} + \frac{1}{PQ} = \frac{1}{r} (\frac{\sin(x + a)}{\sin a} + \frac{\sin(x - b)}{\sin b}) = \frac{\sin(x+a)\sin b + \sin(x-b)\sin a}{r \sin a \sin b} = \frac{\sin x\cos a\sin b + \sin a\cos x\sin b + \sin x\cos b\sin a - \cos x\sin b\sin a}{r \sin a \sin b} = \frac{\sin x(\sin b \cos a + \sin a \cos b)}{r \sin a \sin b}\]

 Clearly, this quantity is maximized when $\sin x = 1.$ Because $x$ must be less than $\pi$, $\frac{1}{PQ} + \frac{1}{PR}$ is as large as possible when $x = \frac{\pi}{2},$ or when line $QR$ is perpendicular to line $PO$.

 


	\item (\textit{USAMO 1979}) Let $A_1,A_2,...,A_{n+1}$ be distinct subsets of $[n]$ with $|A_1|=|A_2|=\cdots =|A_{n+1}|=3$.  Prove that $|A_i\cap A_j|=1$ for some pair $\{i,j\}$. Note that $[n] = \{1, 2, 3, ..., n\}$, or, alternatively, $\{x: 1 \le x \le n\}$.


 



\textbf{Solution}

Suppose the problem statement does not hold. It is clear that $n \ge 4$. Choose the smallest $n$ such that $|A_i \cap A_j| \neq 1$ for each $\{i, j\}$.

 First, the $(n+1)$ subsets have $3(n+1)$ elements (some repeated) in conjunction. Because there are $n$ elements of $[n]$ total, by the Pigeonhole Principle one element of $[n]$, say $k$, is in at least four subsets. Let these subsets be $A_1, A_2, A_3, A_4$, without loss of generality, and let $A_1$ have elements $k, m, n$. Then without loss of generality let $A_2$ have elements $k, m, p$. If set $A_3$ has elements $k, n, p$, then simple casework shows that it is impossible to create $A_4$ without having $A_4$ intersect one of $A_1, A_2, A_3$ at exactly one element. Thus assume $A_3$ has elements $k, m, q$. Then $A_4$ has elements $k, m, r$. Consider each remaining set $A_i$. Then $A_i$ either contains both $k, m$ or none of them. Because there are $(n-2)$ distinct elements of $[n]$ apart from $k, m$, at most $(n-2)$ subsets can contain $k, m$. Then at least 3 subsets do not contain $k, m$, and it is easy to see that they are disjoint from those subsets that do contain $k, m$. Thus, we can partition $[n]$ into two subsets, one of which is the union of the $t$ subsets that do contain $k, m$, and the other is the union of the $(n+1)-t$ subsets that do not contain $k, m$. Because the latter subset has $(n-t-2) < (n+1)-t - 1$ elements, we may use infinite descent to contradict the minimality of $n$. The proof is complete.

 


\end{enumerate}